\section{Finite Externalities and Connected Reconstruction}
\subsection{Direct Pair and Triple Effects}

All 2,016 E03 cases were smooth-valid. Among 672 pair cases, 266 exceeded the
five-budget resolution criterion; among 1,344 triple cases, 147 did so. The
unresolved majority, especially 1,197 of the triple cases, is part of the
result: higher-order effects are not ubiquitous at the frozen amplitudes.

The strongest resolved discovery pair was A7--A8 at OP12,
$\ext{\{7,8\}}\Phi=-0.00535233$ with direct uncertainty
$6.93\times10^{-6}$. The largest locked-holdout value for the same pair was
$-0.00483156$. It combines a 20-MW dispatch increase at GFL37 with a 10\%
reactance increase on its incident branch. The leading triple was A2--A7--A8,
which adds the bus-37 PLL bandwidth action. Its largest discovery coefficient
was $-2.64008\times10^{-5}$, and its largest holdout magnitude was
$2.45036\times10^{-5}$.

The holdout replication rule was met by ten pairs and five triples, each with
same-sign fraction 1.0. (The broader C1 gate counts one additional pair under
its original population rule; the later E05B population contains the ten
fully retained pairs used for anatomy and stress.) These are existential
results on the frozen coordinates. They do not imply that every finite
controller or network portfolio is non-additive.

\subsection{Independent Connected Closure}

All 413 effects resolved in E03 were reconstructed through
\eqref{eq:bridge}--\eqref{eq:connected}. On the locked holdout, all 129 cases
satisfied \eqref{eq:recongate}. Median absolute errors were
$1.11638\times10^{-6}$ for pairs and $5.80806\times10^{-8}$ for triples.
The overall mean bias, $2.45587\times10^{-7}$, was below the median combined
uncertainty. Figure~\ref{fig:reconstruction} shows direct and reconstructed
values with case-specific uncertainty bars; the pointwise gate, rather than
the apparent correlation, supports the closure claim.

\begin{figure}[t]
\centering
\includegraphics[width=\columnwidth]{figure_E04_1_direct_vs_reconstructed.pdf}
\caption{Direct re-equilibrated externality versus connected reconstruction
for every resolved discovery and holdout case. Error bars are conservative
combined uncertainty; the signed symlog axes retain small and negative
values. All 129 holdout cases close within five budgets.}
\label{fig:reconstruction}
\end{figure}

This agreement checks more than a final scalar. The direct route uses only
coalition vertices and M\"obius differencing. The connected route evaluates
total matrix derivatives, the complete set-partition expansion, frequency
integration, and action-hypercube quadrature. Closing the two independently
implemented paths is therefore evidence for the calculus and its numerical
implementation. It is not a causal identification theorem: both calculations
operate on the same model and action definitions.

\subsection{Why Re-Equilibration Is Material to the Calculation}

The connected anatomy distinguishes intrinsic mixed derivatives, products of
single-action feedback terms, and mixed curvature. For A7--A8 on holdout, the
median curvature contribution was $-0.00338136$ and the median feedback
contribution was $-0.000694529$, giving a median singleton-feedback adequacy
ratio of only 0.169. Thus most of this finite pair is associated with the
operating-point/network curvature that appears when dispatch and branch
reactance are changed together.

At OP00, the fully re-equilibrated A7--A8 externality is approximately
$-0.00411$, while the frozen-affine result is $-0.00065321$. The frozen
calculation therefore captures only about 15.9\% of the finite magnitude, an
underestimate of approximately $6.3\times$. Figure~\ref{fig:frozen} shows
the diagnostic for the registered exemplars. This factor belongs to A7--A8 at
the reference operating point; it is not asserted as a universal correction
factor.

\begin{figure}[t]
\centering
\includegraphics[width=\columnwidth]{figure_E04_4_frozen_vs_reequilibrated.pdf}
\caption{Fully re-equilibrated versus frozen-affine externality at OP00 for
registered examples. The dispatch--network pair A7--A8 lies far from the
identity line: re-equilibration increases its magnitude by about $6.3\times$.}
\label{fig:frozen}
\end{figure}

The result sharpens the comparison with local sensitivity. Higher-order
equilibrium sensitivities can be computed
\cite{nam2000eigensensitivity,zamoracardenas2013multiparameter}; the present
calculus does not claim otherwise. Its distinction is that the target object
is the exact finite coalition remainder, and that mixed curvature is retained
along the entire finite action path rather than truncated at one expansion
point.

For the resolved triples, the median adequacy of the pure-feedback
description was 0.975323, but signed terms sometimes canceled strongly.
Accordingly, component signs and magnitudes are reported directly rather than
converted into unstable percentage shares near a zero total. A large
adequacy ratio for one action family does not justify discarding curvature
for another.
